% Options for packages loaded elsewhere
\PassOptionsToPackage{unicode}{hyperref}
\PassOptionsToPackage{hyphens}{url}
\PassOptionsToPackage{dvipsnames,svgnames,x11names}{xcolor}
\documentclass[
]{article}
\usepackage{xcolor}
\usepackage{amsmath,amssymb}
\setcounter{secnumdepth}{-\maxdimen} % remove section numbering
\usepackage{iftex}
\ifPDFTeX
  \usepackage[T1]{fontenc}
  \usepackage[utf8]{inputenc}
  \usepackage{textcomp} % provide euro and other symbols
\else % if luatex or xetex
  \usepackage{unicode-math} % this also loads fontspec
  \defaultfontfeatures{Scale=MatchLowercase}
  \defaultfontfeatures[\rmfamily]{Ligatures=TeX,Scale=1}
\fi
\usepackage{lmodern}
\ifPDFTeX\else
  % xetex/luatex font selection
\fi
% Use upquote if available, for straight quotes in verbatim environments
\IfFileExists{upquote.sty}{\usepackage{upquote}}{}
\IfFileExists{microtype.sty}{% use microtype if available
  \usepackage[]{microtype}
  \UseMicrotypeSet[protrusion]{basicmath} % disable protrusion for tt fonts
}{}
\makeatletter
\@ifundefined{KOMAClassName}{% if non-KOMA class
  \IfFileExists{parskip.sty}{%
    \usepackage{parskip}
  }{% else
    \setlength{\parindent}{0pt}
    \setlength{\parskip}{6pt plus 2pt minus 1pt}}
}{% if KOMA class
  \KOMAoptions{parskip=half}}
\makeatother
\usepackage{longtable,booktabs,array}
\newcounter{none} % for unnumbered tables
\usepackage{calc} % for calculating minipage widths
% Correct order of tables after \paragraph or \subparagraph
\usepackage{etoolbox}
\makeatletter
\patchcmd\longtable{\par}{\if@noskipsec\mbox{}\fi\par}{}{}
\makeatother
% Allow footnotes in longtable head/foot
\IfFileExists{footnotehyper.sty}{\usepackage{footnotehyper}}{\usepackage{footnote}}
\makesavenoteenv{longtable}
\setlength{\emergencystretch}{3em} % prevent overfull lines
\providecommand{\tightlist}{%
  \setlength{\itemsep}{0pt}\setlength{\parskip}{0pt}}
\usepackage{xcolor}
\usepackage{xurl}
\usepackage{unicode-math}
\usepackage{newunicodechar}
\setmainfont{Palatino}
\setmathfont{STIX Two Math}
\setmonofont{Menlo}[Scale=0.9]

% Allow breaks at more characters in URLs
\renewcommand\UrlBreaks{\do\/\do\-\do\.\do\_\do\?\do\&\do\=\do\#\do\:\do\@\do\%}

% Stretch as last resort
\emergencystretch=2em

% Full-width horizontal rules
\let\oldrule\rule
\renewcommand{\rule}[2]{\noindent\oldrule{\linewidth}{0.4pt}}

\newunicodechar{→}{\ensuremath{\rightarrow}}
\newunicodechar{≡}{\ensuremath{\equiv}}
\newunicodechar{≠}{\ensuremath{\neq}}
\newunicodechar{≥}{\ensuremath{\geq}}
\newunicodechar{≤}{\ensuremath{\leq}}
\newunicodechar{∈}{\ensuremath{\in}}
\newunicodechar{∣}{\ensuremath{\mid}}
\newunicodechar{∥}{\ensuremath{\parallel}}
\newunicodechar{ℰ}{\ensuremath{\mathcal{E}}}
\newunicodechar{𝒮}{\ensuremath{\mathcal{S}}}
\newunicodechar{Σ}{\ensuremath{\Sigma}}
\newunicodechar{φ}{\ensuremath{\varphi}}
\newunicodechar{ℤ}{\ensuremath{\mathbb{Z}}}
\usepackage{bookmark}
\IfFileExists{xurl.sty}{\usepackage{xurl}}{} % add URL line breaks if available
\urlstyle{same}
\hypersetup{
  colorlinks=true,
  linkcolor={Maroon},
  filecolor={Maroon},
  citecolor={Blue},
  urlcolor={Blue},
  pdfcreator={LaTeX via pandoc}}

\author{}
\date{}

\begin{document}

\section{Laboratorium nr 10}\label{laboratorium-nr-10}

\subsection{Protokoły dowodzenia z wiedzą zerową --- interaktywny
protokół
Schnorra}\label{protokoux142y-dowodzenia-z-wiedzux105-zerowux105-interaktywny-protokuxf3ux142-schnorra}

\textbf{Liczba punktów:} 3pkt + \textcolor{green}{2pkt}

\textbf{Forma zaliczenia:} live demo + pytania kontrolne

\begin{center}\rule{0.5\linewidth}{0.5pt}\end{center}

\subsection{1. Opis laboratorium}\label{opis-laboratorium}

Celem laboratorium jest poznanie i implementacja interaktywnego
protokołu dowodzenia z wiedzą zerową (\emph{zero-knowledge proof of
knowledge}, ZKPoK) na przykładzie protokołu identyfikacji Schnorra
(Schnorr 1989). W ramach laboratorium:

\begin{enumerate}
\def\labelenumi{\arabic{enumi}.}
\tightlist
\item
  Zaimplementujesz interaktywny protokół Schnorra w grupie cyklicznej
  rzędu pierwszego \(\mathbb{Z}^*_p\);
\item
  Podepniesz go jako klient do serwera autoryzacji udostępnianego przez
  prowadzącego i wykonasz uwierzytelnienie bez przesyłania hasła.
\end{enumerate}

\subsubsection{1.1. Język protokołów
ZKP}\label{jux119zyk-protokoux142uxf3w-zkp}

Protokół dowodzenia z wiedzą zerową jest dwustronnym protokołem między
\textbf{Proverem} \(P\) a \textbf{Verifierem} \(V\). Centralnym pojęciem
jest \textbf{relacja} \(R\) określająca, co znaczy ``dowód poprawny'':

\begin{itemize}
\tightlist
\item
  \(x\) --- \emph{statement} (wypowiedź publiczna, wejście jawne);
\item
  \(w\) --- \emph{witness} (świadek, sekret znany tylko Proverowi);
\item
  \(R(x, w) \in \{0, 1\}\) --- funkcja relacji; Prover twierdzi: ``znam
  \(w\) takie, że \(R(x, w) = 1\)'';
\item
  \(b \in \{0, 1\}\) --- \emph{decision bit} wyjściowy Verifiera (\(1\)
  = accept, \(0\) = reject).
\end{itemize}

Dla protokołu Schnorra relacja jest następująca:

\[R_{\mathrm{DLP}} = \big\{\, ((g, p, q, y),\, x) \;:\; y \equiv g^x \pmod{p} \,\big\}.\]

Statement to czwórka \((g, p, q, y)\), a witness to logarytm dyskretny
\(x\). Prover dowodzi znajomości \(x\) bez ujawniania jego wartości.

\subsubsection{1.2. Trzy własności
bezpieczeństwa}\label{trzy-wux142asnoux15bci-bezpieczeux144stwa}

Protokół ZKP musi spełniać trzy własności, z których każda chroni przed
innym typem ataku i ma odrębną definicję formalną.

\textbf{(1) Completeness (zupełność).} Jeśli \(R(x, w) = 1\) i obie
strony postępują zgodnie z protokołem, to \(V\) wyprowadza \(b = 1\) z
prawdopodobieństwem \(1\).

\textbf{(2) Knowledge Soundness (KSND, niezawodność wiedzy).} Chroni
przed nieuczciwym Proverem \(P^*\). Jeśli \(P^*\) przekonuje \(V\) z
istotnym prawdopodobieństwem, to \textbf{istnieje algorytm Knowledge
Extractor} \(\mathcal{E}\), który mając dostęp do \(P^*\) potrafi
wydobyć poprawnego świadka \(w^*\). Operacyjnie: dwa akceptowane
transkrypty \((a, e_1, z_1)\) i \((a, e_2, z_2)\) z różnymi wyzwaniami
\(e_1 \neq e_2\) pozwalają wyliczyć \(w\). To dowodzi, że \(P^*\)
\emph{musiał} znać sekret.

\textbf{(3) Zero-Knowledge (ZK).} Chroni przed nieuczciwym Verifierem
\(V^*\). \textbf{Istnieje algorytm Simulator} \(\mathcal{S}\), który
mając tylko statement \(x\) (bez witness \(w\)) produkuje transkrypty
\emph{nieodróżnialne} od transkryptów prawdziwej interakcji. To dowodzi,
że \(V^*\) nie uzyskuje żadnej informacji o \(w\), której nie mógłby
uzyskać sam, znając tylko \(x\).

\begin{quote}
\textbf{Uwaga.} Klasyczna definicja ZK zakłada uczciwego Verifiera (HVZK
--- \emph{Honest Verifier Zero-Knowledge}). Pełne ZK wobec dowolnego
\(V^*\) wymaga dodatkowych technik (rewinding, Fischlin transform). W
tym laboratorium pracujemy z HVZK.
\end{quote}

\subsubsection{1.3. Σ-protokoły}\label{ux3c3-protokoux142y}

Protokół Schnorra należy do rodziny \textbf{Σ-protokołów}:
trzy-wiadomościowych protokołów \emph{publicznych monet} (public-coin),
w których wyzwanie \(V\) jest losowane jawnie. Schemat:

\[
\begin{array}{lrcl}
\text{Krok 1 (commitment):} & P \to V & : & a \\
\text{Krok 2 (challenge):}  & V \to P & : & e \\
\text{Krok 3 (response):}   & P \to V & : & z
\end{array}
\]

Verifier akceptuje \textbf{wtedy i tylko wtedy}, gdy
\(\varphi(x, a, e, z) = 1\) dla pewnego predykatu \(\varphi\).

\subsubsection{1.4. Protokół identyfikacji
Schnorra}\label{protokuxf3ux142-identyfikacji-schnorra}

\textbf{Parametry publiczne:} liczba pierwsza \(p\), liczba pierwsza
\(q \mid p-1\), generator \(g\) podgrupy rzędu \(q\) w
\(\mathbb{Z}^*_p\). Zalecane: \(|p| \geq 2048\) b. W laboratorium
używamy parametrów RFC 7919 ffdhe2048 (sekcja 5.1).

\textbf{Statement i witness:} \(x_{\mathrm{pub}} = (g, p, q, y)\);
witness \(w = x_{\mathrm{priv}}\), takie że
\(y = g^{x_{\mathrm{priv}}} \bmod p\).

{\def\LTcaptype{none} % do not increment counter
\begin{longtable}[]{@{}
  >{\raggedright\arraybackslash}p{(\linewidth - 4\tabcolsep) * \real{0.2500}}
  >{\raggedright\arraybackslash}p{(\linewidth - 4\tabcolsep) * \real{0.3333}}
  >{\raggedright\arraybackslash}p{(\linewidth - 4\tabcolsep) * \real{0.4167}}@{}}
\toprule\noalign{}
\begin{minipage}[b]{\linewidth}\raggedright
Krok
\end{minipage} & \begin{minipage}[b]{\linewidth}\raggedright
Strona
\end{minipage} & \begin{minipage}[b]{\linewidth}\raggedright
Operacja
\end{minipage} \\
\midrule\noalign{}
\endhead
\bottomrule\noalign{}
\endlastfoot
1. Commitment & \(P\) & losuje \(r \in_R \{1, \ldots, q-1\}\); oblicza
\(a = g^r \bmod p\); wysyła \(a\) \\
2. Challenge & \(V\) & losuje \(e \in_R \{0, \ldots, 2^t-1\}\); wysyła
\(e\) \\
3. Response & \(P\) & oblicza
\(z = (r + e \cdot x_{\mathrm{priv}}) \bmod q\); wysyła \(z\) \\
4. Verify & \(V\) & akceptuje wtw.
\(g^z \equiv a \cdot y^e \pmod{p}\) \\
\end{longtable}
}

\textbf{Completeness:}

\[g^z = g^{r + e \cdot x_{\mathrm{priv}}} = g^r \cdot \big(g^{x_{\mathrm{priv}}}\big)^e = a \cdot y^e \pmod{p}.\]

\textbf{Knowledge Extractor (intuicja).} Z dwóch akceptowanych
transkryptów \((a, e_1, z_1)\) i \((a, e_2, z_2)\) z różnymi wyzwaniami
\(e_1 \neq e_2\) wynika:

\[z_1 - z_2 \equiv (e_1 - e_2) \cdot x_{\mathrm{priv}} \pmod{q},\]

skąd:

\[x_{\mathrm{priv}} = (z_1 - z_2) \cdot (e_1 - e_2)^{-1} \bmod q.\]

Konstrukcja Extractora dowodzi, że Prover przekonujący Verifiera
dwukrotnie \emph{musi} znać \(x_{\mathrm{priv}}\). Ta sama algebra leży
u podstaw ataku na podpisy Schnorra/ECDSA z powtórzonym nonce (między
innymi incydent Sony PlayStation 3, 2010) --- temat omawiany na
wykładzie nr 10.

\subsubsection{1.5. Zastosowania w systemach
produkcyjnych}\label{zastosowania-w-systemach-produkcyjnych}

Protokół zaimplementowany w niniejszym laboratorium jest dydaktycznym
przedstawieniem matematyki, która stanowi rdzeń następujących systemów
używanych komercyjnie:

\begin{itemize}
\tightlist
\item
  \textbf{Ed25519} (RFC 8032, NIST FIPS 186-5) --- wariant Schnorra na
  krzywej eliptycznej Edwards \(\mathrm{Ed}25519\). Standard SSH, TLS
  1.3, Signal, WhatsApp, WireGuard, Apple Passkeys, Android Keystore.
\item
  \textbf{Bitcoin Taproot} (BIP-340, aktywne od 2021 r.) --- Schnorr na
  krzywej secp256k1, używany do każdej nowej transakcji Bitcoin.
\item
  \textbf{FIDO2 / WebAuthn} --- globalny standard uwierzytelnienia bez
  haseł (Google, Apple, Microsoft); pod spodem ECDSA lub Ed25519.
\item
  \textbf{Konsensus blockchain} --- Solana, Polkadot, Cardano używają
  wariantów Schnorra do podpisywania transakcji w sieciach o
  przepustowości rzędu \(10^4\)--\(10^5\) tps. W każdym z tych systemów
  występują dokładnie te same trzy kroki: commitment, challenge,
  response. Różni je \textbf{grupa}, w której działają, oraz
  \textbf{sposób generowania wyzwania} \(e\). W zapisie multiplikatywnym
  (jak w niniejszym laboratorium) weryfikacja ma postać:
\end{itemize}

\[g^z \equiv a \cdot y^e \pmod{p}.\]

W zapisie addytywnym dla krzywych eliptycznych identyczna zależność
wygląda następująco:

\[z \cdot G \;=\; R \,+\, e \cdot Y,\]

gdzie \(G\) jest punktem bazowym krzywej,
\(Y = x_{\mathrm{priv}} \cdot G\) jest kluczem publicznym,
\(R = r \cdot G\) jest commitment'em, a operacje \(+\) i \(\cdot\) to
odpowiednio dodawanie punktów oraz mnożenie punktu przez skalar.
Algebraicznie jest to ta sama tożsamość, co weryfikacja z sekcji 1.4 ---
różni się tylko notacją podyktowaną wyborem grupy.

Sposób uzyskiwania wyzwania w wersji nieinteraktywnej (podpis cyfrowy)
wykorzystuje transformację Fiat-Shamir:

\[e = H(\mathrm{pk} \,\|\, a \,\|\, m),\]

gdzie \(H\) jest kryptograficzną funkcją skrótu (SHA-256, SHA-512 lub w
przypadku Ed25519 --- SHA-512 z dodatkową strukturą). Pominięcie
\(\mathrm{pk}\), \(a\) lub \(m\) w wejściu \(H\) otwiera atak; temat
jest omawiany na wykładzie nr 10.

\begin{center}\rule{0.5\linewidth}{0.5pt}\end{center}

\subsection{2. Materiały}\label{materiaux142y}

\begin{enumerate}
\def\labelenumi{\arabic{enumi}.}
\tightlist
\item
  C.-P. Schnorr, \emph{Efficient Identification and Signatures for Smart
  Cards}, CRYPTO '89,
  https://link.springer.com/chapter/10.1007/0-387-34805-0\_22
\item
  A. Menezes, P. van Oorschot, S. Vanstone, \emph{Handbook of Applied
  Cryptography}, rozdz. 10,
  http://cacr.uwaterloo.ca/hac/about/chap10.pdf
\item
  D. Boneh, V. Shoup, \emph{A Graduate Course in Applied Cryptography},
  rozdz. 19, https://toc.cryptobook.us/
\item
  A. Fiat, A. Shamir, \emph{How to Prove Yourself}, CRYPTO '86,
  https://link.springer.com/chapter/10.1007/3-540-47721-7\_12
\item
  S. Goldwasser, S. Micali, C. Rackoff, \emph{The Knowledge Complexity
  of Interactive Proof-Systems}, STOC '85,
  https://doi.org/10.1145/22145.22178
\item
  T. Silde, A. Takahashi, \emph{Zero Knowledge Proofs: Challenges,
  Applications, and Real-world Deployment}, NIST Workshop on
  Privacy-Enhancing Cryptography 2024,
  https://csrc.nist.gov/Events/2024/wpec2024
\item
  D. Pointcheval, J. Stern, \emph{Security Arguments for Digital
  Signatures and Blind Signatures}, Journal of Cryptology 2000,
  https://link.springer.com/article/10.1007/s001450010003
\item
  M. Green, \emph{Zero Knowledge Proofs: An Illustrated Primer},
  https://blog.cryptographyengineering.com/2014/11/27/zero-knowledge-proofs-illustrated-primer/
\item
  BIP-340: \emph{Schnorr Signatures for secp256k1},
  https://github.com/bitcoin/bips/blob/master/bip-0340.mediawiki
\item
  RFC 7919: \emph{Negotiated Finite Field Diffie-Hellman Ephemeral
  Parameters for Transport Layer Security},
  https://www.rfc-editor.org/rfc/rfc7919.html
\item
  RFC 8032: \emph{Edwards-Curve Digital Signature Algorithm (EdDSA)},
  https://www.rfc-editor.org/rfc/rfc8032.html
\end{enumerate}

\begin{center}\rule{0.5\linewidth}{0.5pt}\end{center}

\subsection{3. Zadanie do wykonania}\label{zadanie-do-wykonania}

Napisz aplikację realizującą poniższe zadanie. Język programowania
\textbf{dowolny}. Wymagana jest obsługa arytmetyki dużych liczb.

Klient referencyjny w języku Python (\texttt{client.py}) zawiera gotowy
szkielet HTTP, parametry grupy oraz puste funkcje kryptograficzne
(oznaczone \texttt{TODO\ STUDENT}) do uzupełnienia. Plik dostępny jest
na zajęciach.

\subsubsection{\texorpdfstring{3.1. Zadanie 1 (3 pkt +
\textcolor{green}{2 pkt}) --- implementacja interaktywnego protokołu
Schnorra i klienta
uwierzytelniającego}{3.1. Zadanie 1 (3 pkt + ) --- implementacja interaktywnego protokołu Schnorra i klienta uwierzytelniającego}}\label{zadanie-1-3-pkt-implementacja-interaktywnego-protokoux142u-schnorra-i-klienta-uwierzytelniajux105cego}

Zaimplementuj protokół identyfikacji Schnorra zgodnie z opisem w sekcji
1.4 oraz podepnij go jako klienta do serwera pokazowego.

\textbf{Część A: lokalna implementacja.}

\begin{enumerate}
\def\labelenumi{\arabic{enumi}.}
\tightlist
\item
  Pobierz parametry \((p, q, g)\) z serwera (\texttt{GET\ /params}).
  Zweryfikuj lokalnie:

  \begin{itemize}
  \tightlist
  \item
    \(g^q \equiv 1 \pmod{p}\) (generator ma rząd dzielący \(q\));
  \item
    \(g \neq 1\).
  \end{itemize}
\item
  Zaimplementuj funkcję \texttt{keygen()} zwracającą parę
  \((x_{\mathrm{priv}}, y)\), gdzie
  \(x_{\mathrm{priv}} \in_R \{1, \ldots, q-1\}\) oraz
  \(y = g^{x_{\mathrm{priv}}} \bmod p\).
\item
  Zaimplementuj cztery rozdzielne funkcje protokołu:

  \begin{itemize}
  \tightlist
  \item
    \texttt{prover\_commit()\ →\ (r,\ a)};
  \item
    \texttt{verifier\_challenge(t\_bits)\ →\ e};
  \item
    \texttt{prover\_response(r,\ e,\ x\_priv)\ →\ z};
  \item
    \texttt{verifier\_check(y,\ a,\ e,\ z)\ →\ bool}.
  \end{itemize}
\item
  Wykonaj testy lokalne:

  \begin{itemize}
  \tightlist
  \item
    \textbf{test completeness:} 100 przebiegów z poprawnym kluczem ---
    wszystkie akceptowane;
  \item
    \textbf{test odporności:} mutacja \(z\) lub \(a\) o \(1\) →
    odrzucenie;
  \item
    \textbf{test z fałszywym kluczem:} Prover używa
    \(x' \neq x_{\mathrm{priv}}\) → odrzucenie.
  \end{itemize}
\end{enumerate}

\textbf{Część B: klient uwierzytelniający.}

\begin{enumerate}
\def\labelenumi{\arabic{enumi}.}
\tightlist
\item
  Zaimplementuj \texttt{register(username,\ y)} --- wywołanie
  \texttt{POST\ /v1/register}.
\item
  Zaimplementuj \texttt{login(username,\ x\_priv)} jako pełną sekwencję:

  \begin{itemize}
  \tightlist
  \item
    obliczenie \((r, a)\) lokalnie przez \texttt{prover\_commit()};
  \item
    \texttt{POST\ /v1/login/start\ \{username,\ a\}} --- otrzymanie
    wyzwania \(e\) oraz \texttt{session\_id};
  \item
    obliczenie \(z = (r + e \cdot x_{\mathrm{priv}}) \bmod q\) lokalnie
    przez \texttt{prover\_response};
  \item
    \texttt{POST\ /v1/login/finish\ \{session\_id,\ z\}} --- otrzymanie
    tokenu sesji (status \(200\)) lub odrzucenia (\(401\)).
  \end{itemize}
\end{enumerate}

\textbf{Demonstracja do oceny.} Live demo przed prowadzącym składa się z
trzech kroków wykonanych na żywo na uruchomionym kliencie:

\begin{enumerate}
\def\labelenumi{\arabic{enumi}.}
\tightlist
\item
  Rejestracja konta studenta (np. \texttt{nazwisko@lab}) --- odpowiedź
  \(200\);
\item
  Udane uwierzytelnienie --- odpowiedź \(200\) z tokenem sesji;
\item
  Próba uwierzytelnienia z podstawioną fałszywą wartością \(z\) (np.
  zwiększoną o \(1\)) --- odpowiedź \(401\).
\end{enumerate}

\begin{center}\rule{0.5\linewidth}{0.5pt}\end{center}

\subsection{4. Pytania kontrolne}\label{pytania-kontrolne}

\begin{enumerate}
\def\labelenumi{\arabic{enumi}.}
\tightlist
\item
  Dlaczego wyzwanie \(e\) musi być losowane \emph{po} wysłaniu
  commitment \(a\)? Jakie konsekwencje miałoby odwrócenie kolejności?
\item
  Klucz prywatny \(x_{\mathrm{priv}}\) nigdy nie opuszcza klienta ---
  serwer otrzymuje tylko \(y\), \(a\) i \(z\). Jaką to ma przewagę nad
  uwierzytelnieniem hasłem w sytuacji, gdy serwer zostaje
  skompromitowany (np. wyciekła baza danych)?
\end{enumerate}

\begin{center}\rule{0.5\linewidth}{0.5pt}\end{center}

\subsection{5. Wskazówki
implementacyjne}\label{wskazuxf3wki-implementacyjne}

\subsubsection{5.1. Parametry grupy (RFC 7919
ffdhe2048)}\label{parametry-grupy-rfc-7919-ffdhe2048}

Identyczne wartości używa serwer pokazowy. Pełna postać liczbowa w pliku
\texttt{client.py}. Schematycznie:

\[
\begin{aligned}
p &= \mathtt{FFFFFFFF}\ldots\mathtt{FFFFFFFF} \quad (2048 \text{ bitów}) \\
q &= (p - 1)/2 \quad (2047 \text{ bitów},\ \textit{safe prime}) \\
g &= 2
\end{aligned}
\]

Wymagana lokalna weryfikacja: \(g^q \equiv 1 \pmod{p}\) oraz
\(g \neq 1\).

\subsubsection{5.2. Serializacja liczb przy komunikacji z
serwerem}\label{serializacja-liczb-przy-komunikacji-z-serwerem}

Wszystkie wartości liczbowe przesyłane są w formacie \textbf{lowercase
hex string} (bez prefiksu \texttt{0x}). Konwencja jednolita dla obu
kierunków komunikacji. Format JSON Number nie jest stosowany z uwagi na
\(53\)-bitowe ograniczenie precyzji.

\subsubsection{5.3. Losowanie liczb}\label{losowanie-liczb}

Do losowania liczb stosuj wyłącznie kryptograficznie bezpieczne
generatory:

\begin{itemize}
\tightlist
\item
  Python \(\geq 3.8\): moduł \texttt{secrets}, funkcja
  \texttt{secrets.randbelow(q\ -\ 1)\ +\ 1};
\item
  Java: \texttt{java.security.SecureRandom};
\item
  C/C++: \texttt{/dev/urandom}, biblioteka GMP z dobrym źródłem
  entropii;
\item
  Rust: crate \texttt{rand} z \texttt{OsRng}.
\end{itemize}

Nie należy używać generatorów pseudolosowych przeznaczonych do
zastosowań niekryptograficznych (\texttt{random.randint} w Pythonie,
\texttt{java.util.Random}, \texttt{rand()} w C).

\subsubsection{5.4. Arytmetyka modularna}\label{arytmetyka-modularna}

Najczęstszym źródłem błędów w implementacji protokołu jest pomylenie
modulo. Obowiązuje zasada:

\begin{itemize}
\tightlist
\item
  bazę redukujemy \(\bmod\ p\) (operacje na elementach grupy
  \(\mathbb{Z}^*_p\));
\item
  wykładniki redukujemy \(\bmod\ q\) (operacje w pierścieniu
  \(\mathbb{Z}_q\), rząd podgrupy).
\end{itemize}

W szczególności: \(z = (r + e \cdot x_{\mathrm{priv}}) \bmod q\),
\textbf{nie} \(\bmod\ p\).

\subsubsection{5.5. Długość
wyzwania}\label{dux142ugoux15bux107-wyzwania}

W laboratorium przyjmujemy \(t = 256\) bitów, tj.
\(e \in_R \{0, \ldots, 2^{256}-1\}\). Wartość zgodna z długością wyjścia
SHA-256. Błąd \emph{knowledge soundness} wynosi \(2^{-256}\) na sesję.
Wartość ta jest stała i nie jest udostępniana przez \texttt{/params} ---
klient zakłada ją z góry.

\subsubsection{5.6. Endpointy serwera
pokazowego}\label{endpointy-serwera-pokazowego}

{\def\LTcaptype{none} % do not increment counter
\begin{longtable}[]{@{}
  >{\raggedright\arraybackslash}p{(\linewidth - 4\tabcolsep) * \real{0.4348}}
  >{\raggedright\arraybackslash}p{(\linewidth - 4\tabcolsep) * \real{0.3478}}
  >{\raggedright\arraybackslash}p{(\linewidth - 4\tabcolsep) * \real{0.2174}}@{}}
\toprule\noalign{}
\begin{minipage}[b]{\linewidth}\raggedright
Endpoint
\end{minipage} & \begin{minipage}[b]{\linewidth}\raggedright
Metoda
\end{minipage} & \begin{minipage}[b]{\linewidth}\raggedright
Cel
\end{minipage} \\
\midrule\noalign{}
\endhead
\bottomrule\noalign{}
\endlastfoot
\texttt{/params} & GET & parametry grupy \((p, q, g)\) \\
\texttt{/v1/register} & POST & rejestracja konta z kluczem publicznym
\(y\) \\
\texttt{/v1/login/start} & POST & rozpoczęcie uwierzytelnienia
(commitment \(a\)) \\
\texttt{/v1/login/finish} & POST & zakończenie uwierzytelnienia
(response \(z\)) \\
\end{longtable}
}

Adres serwera podaje prowadzący na zajęciach.

\begin{center}\rule{0.5\linewidth}{0.5pt}\end{center}

\textbf{Termin oddania zadania:} do laboratorium nr 12 włącznie. Demo po
tym terminie nie jest punktowane.

\end{document}
